% !TEX root = 概率论期中复习.tex
% 自动抽取自 chapters/chapter1.tex；原例题及其解答保留。

\clearpage\section{随机试验和事件}\label{k:events}
事件的全集为$\Omega$；$AB=A\cap B$，$A^c=\Omega\setminus A$。并表示至少一个发生，交表示同时发生，差表示前者发生而后者不发生。

\begin{enumerate}
\item \textbf{互斥（互不相容）事件}

若 $A\cap B=\varnothing$，则称事件 $A$ 与 $B$ 互斥或互不相容，其概率意义是二者不能同时发生。若对一列事件 $A_n\ (n\geq1)$，任意 $i\ne j$ 都有 $A_i\cap A_j=\varnothing$，则称该列事件两两互斥。

\item \textbf{补事件}

事件 $A$ 的补集
\[
A^c=\left\{\omega\in\Omega\mathrel{\big|}\omega\notin A\right\}
\]
称为 $A$ 的补事件，其概率意义是 $A$ 不发生。并且
\[
A\setminus B=A\cap B^c.
\]
事件运算约定先进行补运算，再进行交运算，最后进行并运算或差运算。

\item \textbf{完备事件组}

若事件 $A_1,A_2,\ldots$ 两两互斥，且
\[
\bigcup_{i=1}^{\infty}A_i=\Omega,
\]
则称 $A_1,A_2,\ldots$ 构成一个\textbf{完备事件组}，或称为 $\Omega$ 的一个可数划分。
\end{enumerate}

事件的运算性质就是集合的运算性质。对任意事件 $A,B,C$，有
\[
A\cup B=B\cup A,\qquad A\cap B=B\cap A,
\]
\[
(A\cup B)\cup C=A\cup(B\cup C),\qquad
(A\cap B)\cap C=A\cap(B\cap C),
\]
\[
A\cup(B\cap C)=(A\cup B)\cap(A\cup C),\qquad
A\cap(B\cup C)=(A\cap B)\cup(A\cap C).
\]
此外，德摩根律为
\[
\left(\bigcup_{i=1}^{n}A_i\right)^c=\bigcap_{i=1}^{n}A_i^c,
\qquad
\left(\bigcap_{i=1}^{n}A_i\right)^c=\bigcup_{i=1}^{n}A_i^c,
\]
以及
\[
\left(\bigcup_{i=1}^{\infty}A_i\right)^c=\bigcap_{i=1}^{\infty}A_i^c,
\qquad
\left(\bigcap_{i=1}^{\infty}A_i\right)^c=\bigcup_{i=1}^{\infty}A_i^c.
\]

\begin{mydefinition}{随机试验}{random-experiment}
在概率论中，将满足下列条件的试验称为随机试验，简称试验：
\begin{enumerate}
\item 试验进行前，无法确切知道一定会出现哪个结果；
\item 已知试验所有可能的结果；
\item 试验结束后，能够明确知道哪个试验结果出现。
\end{enumerate}
\end{mydefinition}
\origin{原笔记第1.1节，定义“随机试验”；源码第12行。}

\begin{mydefinition}{样本点与样本空间}{sample-space}
\begin{enumerate}
\item 一个随机试验的任意一个可能结果称为一个\textbf{样本点}（sample），记为 $\omega$；
\item 该随机试验所有可能结果组成的集合称为该试验的\textbf{样本空间}（sample space），记为 $\Omega$。
\end{enumerate}
\end{mydefinition}
\origin{原笔记第1.1节，定义“样本点与样本空间”；源码第21行。}

\begin{mydefinition}{事件}{event}
\begin{enumerate}
\item 样本空间 $\Omega$ 的一个子集称为一个\textbf{事件}，通常用大写字母 $A,B,C,\ldots$ 表示；
\item 对事件 $A$，若试验结果 $\omega\in A$，则称事件 $A$ 发生；否则称事件 $A$ 不发生。
\end{enumerate}
\end{mydefinition}
\origin{原笔记第1.1节，定义“事件”；源码第34行。}

\begin{mydefinition}{古典概率模型}{classical-probability-model}
若某随机试验的样本空间 $\Omega$ 只含有限个样本点，且每个样本点出现的概率相等，则称该随机试验为\textbf{古典概率模型}，简称\textbf{古典概型}。对任意事件 $A\subseteq\Omega$，记 $A$ 中元素个数为 $n(A)$，则
\[
P(A)=\dfrac{n(A)}{n(\Omega)}.
\]
\end{mydefinition}
\origin{原笔记第1.1节，定义“古典概率模型”；源码第142行。}

\begin{proposition}{计数的乘法原理}{multiplication-principle}
若一个过程由 $r$ 个试验构成，第 $i$ 个试验有 $n_i$ 种可能结果，则该过程共有
\[
n_1n_2\cdots n_r
\]
种可能结果。
\end{proposition}
\origin{原笔记第1.1节，命题“计数的乘法原理”；源码第151行。}

\subsection{例题：抽取方式与排列计数}\label{note:c1:e1}
\paragraph{题目}
设 $n,m$ 为正整数。求下列计数问题中的不同结果数：
\begin{enumerate}
\item 从 $n$ 个不同的球中有放回地每次取一个，依次取 $m$ 次，按取出顺序排成一列。
\item 从 $n$ 个不同的球中无放回地取出 $m$ 个球，按顺序排成一列，其中 $m\leq n$。
\item $n$ 个不同的人围圆桌坐一圈，将旋转后重合的坐法视为同一种坐法。
\item 从 $n$ 个不同的球中无放回地取出 $m$ 个球，不计顺序，其中 $m\leq n$。
\end{enumerate}

\basis{原笔记第1.1节，例“抽取方式与排列计数”；源码\texttt{chapters/chapter1.tex}第158行。}
\begin{proof}[解答]
\begin{enumerate}
\item 每次都有 $n$ 种选择，由乘法原理，结果数为
\[n^m.\]
\item 依次有 $n,n-1,\ldots,n-m+1$ 种选择，因此排列数为
\[A_n^m=n(n-1)\cdots(n-m+1)=\frac{n!}{(n-m)!}.\]
\item 固定其中一人的位置，其余 $n-1$ 人按顺序安排，恰好消除旋转造成的重复，故圆排列数为
\[(n-1)!.\]
\item 每个不计顺序的组合对应 $m!$ 个有序排列，因此组合数为
\[C_n^m=\binom nm=\frac{A_n^m}{m!}=\frac{n!}{(n-m)!m!}.\]
\end{enumerate}
\end{proof}

\begin{proposition}{多组分配}{multinomial-count}
将 $n$ 个不同元素分成 $r$ 个组，其中第 $i$ 组有 $n_i$ 个元素，组内不计顺序，且
\[
n_1+n_2+\cdots+n_r=n,
\]
则不同分法共有
\[
\dfrac{n!}{n_1!n_2!\cdots n_r!}
\]
种。
\end{proposition}
\origin{原笔记第1.1节，命题“多组分配”；源码第175行。}

\begin{mycorollary}{多项式定理}{multinomial-theorem}
\[
(x_1+x_2+\cdots+x_r)^n
=\sum_{n_1+n_2+\cdots+n_r=n}
\dfrac{n!}{n_1!n_2!\cdots n_r!}
x_1^{n_1}x_2^{n_2}\cdots x_r^{n_r}.
\]
其中
\[
\binom{n}{n_1,n_2,\ldots,n_r}
:=\dfrac{n!}{n_1!n_2!\cdots n_r!}
\]
称为多项式系数。
\end{mycorollary}
\origin{原笔记第1.1节，推论“多项式定理”；源码第202行。}

\begin{proposition}{正整数解的个数}{positive-integer-solutions}
设 $n\geq1$，方程
\[
x_1+x_2+\cdots+x_r=n,\qquad x_i>0\quad(i=1,2,\ldots,r)
\]
有
\[
\binom{n-1}{r-1}
\]
个不同的正整数解。
\end{proposition}
\origin{原笔记第1.1节，命题“正整数解的个数”；源码第217行。}

\begin{proposition}{非负整数解的个数}{nonnegative-integer-solutions}
设 $n\geq1$，方程
\[
x_1+x_2+\cdots+x_r=n,\qquad x_i\geq0\quad(i=1,2,\ldots,r)
\]
有
\[
\binom{n+r-1}{r-1}
\]
个不同的非负整数解。
\end{proposition}
\origin{原笔记第1.1节，命题“非负整数解的个数”；源码第238行。}

\begin{mydefinition}{几何概率模型}{geometric-probability-model}
设随机试验的样本空间 $\Omega$ 是 $\mathbb{R}^d$ 中一个具有体积 $m(\Omega)$ 的子集，且
\[
0<m(\Omega)<+\infty.
\]
若事件 $A\subseteq\Omega$ 具有体积 $m(A)$，并规定
\[
P(A)=\dfrac{m(A)}{m(\Omega)},
\]
则称该随机试验为\textbf{几何概率模型}，简称\textbf{几何概型}。
\end{mydefinition}
\origin{原笔记第1.1节，定义“几何概率模型”；源码第263行。}

\subsection{例题：贝特朗奇论}\label{note:c1:e2}
\paragraph{题目}
随机地从半径为 $r$ 的圆中取一根弦，求该弦长度大于圆内接正三角形边长 $\sqrt{3}r$ 的概率。设事件 $A$ 表示“弦长大于 $\sqrt{3}r$”。

\begin{enumerate}
\item 以弦中点 $M$ 确定弦，规定 $M$ 在圆内均匀取值，求 $P(A)$。
\item 固定弦的一个端点，规定弦与该端点处切线的夹角 $\theta$ 在 $[0,\pi]$ 上均匀取值，求 $P(A)$。
\item 规定弦到圆心的距离 $D$ 在 $[0,r]$ 上均匀取值，求 $P(A)$。
\end{enumerate}
说明三种结果为何不同。

\basis{原笔记第1.1节，例“贝特朗奇论”；源码\texttt{chapters/chapter1.tex}第279行。}
\begin{proof}[解答]
“随机取弦”的说法本身并未说明等可能的参照标准，因此可以得到不同的概率。
\begin{enumerate}
\item 若以弦的中点 $M$ 确定一根弦，并规定 $M$ 在大圆内均匀取值，则弦长大于 $\sqrt{3}r$ 当且仅当 $M$ 落在半径为 $r/2$ 的同心小圆内。因此
\[
P(A)=\dfrac{\pi(r/2)^2}{\pi r^2}=\dfrac14.
\]

\item 若固定弦的一个端点，并规定弦与该端点处切线的夹角 $\theta$ 在 $[0,\pi]$ 上均匀取值，则
\[
A=\left(\dfrac{\pi}{3},\dfrac{2\pi}{3}\right),
\qquad
P(A)=\dfrac{\pi/3}{\pi}=\dfrac13.
\]

\item 若以弦到圆心的距离 $D$ 确定一根弦，并规定 $D$ 在 $[0,r]$ 上均匀取值，则 $A=\{0\leq D<r/2\}$，从而
\[
P(A)=\dfrac{r/2}{r}=\dfrac12.
\]
\end{enumerate}

三种理解得到不同答案的原因在于：它们分别选择了不同的样本空间和等可能标准。求概率前，必须先明确随机试验的样本空间以及概率模型。
\end{proof}

\subsection{例题：布丰投针实验}\label{note:c1:e3}
\paragraph{题目}
在平面上画出间距为 $d$ 的一组平行直线。随机投掷一根长度为 $l$ 的针，设
\[
0<l\leq d.
\]
记针与任一直线相交为事件 $A$。令 $X$ 表示针的中点到最近平行线的距离，$\Theta$ 表示针与平行线的夹角。按照几何概型，取
\[
0\leq X\leq \dfrac d2,
\qquad
0\leq \Theta\leq \dfrac{\pi}{2},
\]
并把 $(X,\Theta)$ 看作在该矩形内等可能取值。

\begin{enumerate}
\item 求针与任一平行线相交的概率 $P(A)$。
\item 若独立投针 $N$ 次，其中 $H>0$ 次与直线相交，给出 $N$ 较大时 $\pi$ 的实验估计式。
\end{enumerate}

\basis{原笔记第1.1节，例“布丰投针实验”；源码\texttt{chapters/chapter1.tex}第305行。}
\begin{proof}[解答]
\begin{enumerate}
\item 针与直线相交当且仅当
\[
X\leq \dfrac l2\sin\Theta.
\]
因此，事件 $A$ 对应的面积为
\[
\int_0^{\pi/2}\dfrac l2\sin\theta\,\mathrm{d}\theta
=\dfrac l2,
\]
而样本空间的面积为 $\dfrac d2\cdot\dfrac{\pi}{2}=\dfrac{d\pi}{4}$。故
\[
P(A)=\dfrac{l/2}{d\pi/4}=\dfrac{2l}{\pi d}.
\]
\item 若独立投针 $N$ 次，其中有 $H$ 次与直线相交，则当 $N$ 较大时，频率 $H/N$ 可近似看作 $P(A)$，从而得到 $\pi$ 的实验估计式
\[
\pi\approx\dfrac{2lN}{dH}.
\]
\end{enumerate}
\end{proof}

\clearpage\section{可测空间和概率空间}\label{k:probability}
事件列$A_n\uparrow A$表示$A_n\subseteq A_{n+1}$且$A=\bigcup_nA_n$；$A_n\downarrow A$表示$A_{n+1}\subseteq A_n$且$A=\bigcap_nA_n$。

\begin{mydefinition}{频率}{frequency}
在 $N$ 次重复试验中，若事件 $A$ 出现 $N(A)$ 次，则比值
\[
f_N(A)=\dfrac{N(A)}{N}
\]
称为事件 $A$ 发生的\textbf{频率}。
\end{mydefinition}
\origin{原笔记第1.2节，定义“频率”；源码第363行。}

\begin{proposition}{频率的性质}{frequency-properties}
对任意事件 $A$，有
\[
0\leq f_N(A)\leq1,
\qquad
f_N(\Omega)=1,
\qquad
f_N(\varnothing)=0.
\]
并且
\[
f_N(A^c)=1-f_N(A).
\]
若 $A_1,A_2,\ldots,A_r$ 两两互斥，则
\[
f_N\left(\bigcup_{i=1}^{r}A_i\right)
=\sum_{i=1}^{r}f_N(A_i).
\]
\end{proposition}
\origin{原笔记第1.2节，命题“频率的性质”；源码第371行。}

\begin{mydefinition}{$\sigma$-域与可测空间}{sigma-algebra}
设 $\Omega$ 是非空集合，$\mathscr{F}$ 是由 $\Omega$ 的一些子集构成的集合类。若 $\mathscr{F}$ 满足：
\begin{enumerate}
\item $\Omega\in\mathscr{F}$；
\item 若 $A\in\mathscr{F}$，则 $A^c\in\mathscr{F}$；
\item 若 $A_n\in\mathscr{F}$（$n\geq1$），则
\[
\bigcup_{n=1}^{\infty}A_n\in\mathscr{F},
\]
\end{enumerate}
则称 $\mathscr{F}$ 为 $\Omega$ 的一个\textbf{$\sigma$-域}或\textbf{事件域}，在分析学中我们称之为\textbf{$\sigma$-代数}，$\mathscr{F}$ 中的元素称为事件或可测集，二元组 $(\Omega,\mathscr{F})$ 称为\textbf{可测空间}。
\end{mydefinition}
\origin{原笔记第1.2节，定义“$\sigma$-域与可测空间”；源码第412行。}

\begin{proposition}{$\sigma$-域的基本封闭性}{sigma-algebra-closure}
若 $\mathscr{F}$ 是 $\Omega$ 的一个 $\sigma$-域，则：
\begin{enumerate}
\item $\varnothing\in\mathscr{F}$；
\item 若 $A_n\in\mathscr{F}$（$n\geq1$），则
\[
\bigcap_{n=1}^{\infty}A_n\in\mathscr{F};
\]
\item 若 $A_1,\ldots,A_k\in\mathscr{F}$，则有限并与有限交都属于 $\mathscr{F}$；
\item 若 $A,B\in\mathscr{F}$，则 $A\setminus B\in\mathscr{F}$。
\end{enumerate}
\end{proposition}
\origin{原笔记第1.2节，命题“$\sigma$-域的基本封闭性”；源码第425行。}

\subsection{例题：\texorpdfstring{一些 $\sigma$-代数的例子}{一些 sigma-代数的例子}}\label{note:c1:e4}
\paragraph{题目}
设 $\Omega$ 为非空集合，$A\subseteq\Omega$。验证下列集合类是 $\sigma$-域，并验证前两项所述的最小性：
\begin{enumerate}
\item $\{\varnothing,\Omega\}$ 是 $\Omega$ 上所含事件最少的 $\sigma$-域；
\item 对任意 $A\subseteq\Omega$，集合类
\[
\{\varnothing,\Omega,A,A^c\}
\]
是包含 $A$ 的最小的 $\sigma$-域；
\item $\mathcal{P}(\Omega):=\{A\mid A\subseteq\Omega\}$ 是 $\Omega$ 的所有子集构成的 $\sigma$-域。
\end{enumerate}

\basis{原笔记第1.2节，例“一些 $\sigma$-代数的例子”；源码\texttt{chapters/chapter1.tex}第451行。}
\begin{proof}[验证]
\begin{enumerate}
\item 集合类 $\{\varnothing,\Omega\}$ 含 $\Omega$，取补仍在其中；其中任意可列个集合的并只能是 $\varnothing$ 或 $\Omega$，所以它是 $\sigma$-域。任意 $\sigma$-域都必须含 $\Omega$ 及其补集 $\varnothing$，故它所含事件最少。
\item 将 $\Omega$ 分为 $A$ 与 $A^c$ 两部分，所列集合恰为这两部分的所有可能并集，故取补和可列并仍为其中一个集合。它含 $\Omega$，所以是 $\sigma$-域。任意包含 $A$ 的 $\sigma$-域还必须包含 $A^c,\varnothing,\Omega$，故所列集合类是包含 $A$ 的最小 $\sigma$-域。若 $A$ 为空集或全集，重复项合并后即为第一项。
\item $\mathcal P(\Omega)$ 含 $\Omega$，其任意元素的补集和任意可列个元素的并仍是 $\Omega$ 的子集，故仍属于 $\mathcal P(\Omega)$。因此它是 $\sigma$-域。
\end{enumerate}
\end{proof}

\begin{proposition}{$\sigma$-域的交}{intersection-of-sigma-algebras}
设 $I$ 是一个指标集。若对任意 $\alpha\in I$，$\mathscr{F}_\alpha$ 都是 $\Omega$ 上的 $\sigma$-域，则
\[
\mathscr{F}=\bigcap_{\alpha\in I}\mathscr{F}_\alpha
\]
也是 $\Omega$ 上的 $\sigma$-域。
\end{proposition}
\origin{原笔记第1.2节，命题“$\sigma$-域的交”；源码第465行。}

\begin{mytheorem}{由集合类生成的 $\sigma$-域}{generated-sigma-algebra}
设 $\Omega$ 是非空集合，$\mathscr{E}$ 是由 $\Omega$ 的一些子集构成的非空集合类。则存在唯一的 $\sigma$-域 $\mathscr{F}$ 满足：
\begin{enumerate}
\item $\mathscr{E}\subseteq\mathscr{F}$；
\item 对任意包含 $\mathscr{E}$ 的 $\sigma$-域 $\mathcal{M}$，都有 $\mathscr{F}\subseteq\mathcal{M}$。
\end{enumerate}
此时称 $\mathscr{F}$ 为由 $\mathscr{E}$ \textbf{生成的 $\sigma$-域}，记为 $\sigma(\mathscr{E})$。
\end{mytheorem}
\origin{原笔记第1.2节，定理“由集合类生成的 $\sigma$-域”；源码第477行。}

\begin{mydefinition}{Borel $\sigma$-域}{borel-sigma-algebra}
记
\[
\mathscr{E}_{\mathbb{R}}=\{(a,b]\mid a,b\in\mathbb{R}\}.
\]
称 $\sigma(\mathscr{E}_{\mathbb{R}})$ 为 $\mathbb{R}$ 上的\textbf{Borel $\sigma$-域}，记为 $\mathscr{B}(\mathbb{R})$。

类似地，记
\[
\mathscr{E}_{\mathbb{R}^d}
=\left\{(a_1,b_1]\times\cdots\times(a_d,b_d]\mid a_i,b_i\in\mathbb{R},\ 1\leq i\leq d\right\}.
\]
称 $\sigma(\mathscr{E}_{\mathbb{R}^d})$ 为 $\mathbb{R}^d$ 上的 Borel $\sigma$-域，记为 $\mathscr{B}(\mathbb{R}^d)$。
\end{mydefinition}
\origin{原笔记第1.2节，定义“Borel $\sigma$-域”；源码第495行。}

\begin{proposition}{常用集合属于 $\mathscr{B}(\mathbb{R})$}{common-borel-sets}
对任意 $a<b$，开区间 $(a,b)$、闭区间 $[a,b]$、半开半闭区间 $[a,b)$ 以及单点集 $\{b\}$ 都属于 $\mathscr{B}(\mathbb{R})$。
\end{proposition}
\origin{原笔记第1.2节，命题“常用集合属于 $\mathscr{B}(\mathbb{R})$”；源码第510行。}

\begin{mydefinition}{概率与概率空间}{probability-space}
设 $\Omega$ 是非空集合，$\mathscr{F}$ 是 $\Omega$ 上给定的一个 $\sigma$-域。若映射
\[
P:\mathscr{F}\longrightarrow[0,+\infty)
\]
满足：
\begin{enumerate}
\item \textbf{非负性：}对任意 $A\in\mathscr{F}$，有 $P(A)\geq0$；
\item \textbf{规范性：}$P(\Omega)=1$；
\item \textbf{可列可加性：}若 $A_n\in\mathscr{F}$（$n\geq1$）两两互斥，则
\[
P\left(\bigcup_{n=1}^{\infty}A_n\right)
=\sum_{n=1}^{\infty}P(A_n),
\]
\end{enumerate}
则称 $P$ 为 $\mathscr{F}$ 上的\textbf{概率}，三元组 $(\Omega,\mathscr{F},P)$ 称为\textbf{概率空间}。
\end{mydefinition}
\origin{原笔记第1.2节，定义“概率与概率空间”；源码第536行。}

\begin{proposition}{空集的概率}{probability-empty-set}
\[
P(\varnothing)=0.
\]
\end{proposition}
\origin{原笔记第1.2节，命题“空集的概率”；源码第555行。}

\begin{proposition}{有限可加性}{finite-additivity}
若 $A_1,\ldots,A_n\in\mathscr{F}$ 两两互斥，则
\[
P\left(\bigcup_{i=1}^{n}A_i\right)
=\sum_{i=1}^{n}P(A_i).
\]
\end{proposition}
\origin{原笔记第1.2节，命题“有限可加性”；源码第576行。}

\begin{proposition}{单调性}{probability-monotonicity}
若 $A,B\in\mathscr{F}$ 且 $B\subseteq A$，则
\[
P(B)\leq P(A).
\]
\end{proposition}
\origin{原笔记第1.2节，命题“单调性”；源码第588行。}

\begin{mycorollary}{概率的上界}{probability-upper-bound}
对任意 $A\in\mathscr{F}$，有
\[
0\leq P(A)\leq1.
\]
\end{mycorollary}
\origin{原笔记第1.2节，推论“概率的上界”；源码第602行。}

\begin{proposition}{补事件的概率}{probability-complement}
对任意 $A\in\mathscr{F}$，有
\[
P(A^c)=1-P(A).
\]
\end{proposition}
\origin{原笔记第1.2节，命题“补事件的概率”；源码第613行。}

\begin{proposition}{一般加法公式}{probability-addition-formula}
对任意事件 $A,B\in\mathscr{F}$，有
\[
P(A\cup B)=P(A)+P(B)-P(A\cap B).
\]
\end{proposition}
\origin{原笔记第1.2节，命题“一般加法公式”；源码第627行。}

\begin{mytheorem}{容斥恒等式}{inclusion-exclusion}
对任意 $n$ 个事件 $A_1,A_2,\ldots,A_n$，有
\[
P\left(\bigcup_{i=1}^{n}A_i\right)
=\sum_{r=1}^{n}(-1)^{r+1}
\sum_{1\leq i_1<\cdots<i_r\leq n}
P\left(A_{i_1}\cap\cdots\cap A_{i_r}\right).
\]
\end{mytheorem}
\origin{原笔记第1.2节，定理“容斥恒等式”；源码第655行。}

\begin{mycorollary}{有限次次可加性}{finite-subadditivity}
对任意 $n$ 个事件 $A_1,\ldots,A_n$，有
\[
P\left(\bigcup_{i=1}^{n}A_i\right)
\leq\sum_{i=1}^{n}P(A_i).
\]
\end{mycorollary}
\origin{原笔记第1.2节，推论“有限次次可加性”；源码第677行。}

\subsection{例题：Kounias 不等式}\label{note:c1:e5}
\paragraph{题目}
证明：对任意 $n$ 个事件 $A_1,\ldots,A_n$，有
\[
P\left(\bigcup_{r=1}^{n}A_r\right)
\leq\min_{1\leq k\leq n}
\left\{
\sum_{r=1}^{n}P(A_r)-\sum_{\substack{1\leq r\leq n\\ r\ne k}}P(A_r\cap A_k)
\right\}.
\]

\basis{原笔记第1.2节，例“Kounias 不等式”；源码\texttt{chapters/chapter1.tex}第696行。}
\begin{proof}
固定 $k$。由有限次次可加性，
\[
\begin{aligned}
\sum_{r=1}^{n}P(A_r)-\sum_{r\ne k}P(A_r\cap A_k)
&=P(A_k)+\sum_{r\ne k}P(A_r\cap A_k^c)\\
&\geq P\left(A_k\cup\bigcup_{r\ne k}(A_r\cap A_k^c)\right)\\
&=P\left(\bigcup_{r=1}^{n}A_r\right).
\end{aligned}
\]
对 $k=1,\ldots,n$ 取最小值即得结论。
\end{proof}

\begin{mydefinition}{事件序列的极限}{limit-of-event-sequences}
若事件列 $\{A_n\}_{n\geq1}$ 满足
\[
A_1\subseteq A_2\subseteq\cdots\subseteq A_n\subseteq\cdots,
\]
则称其为\textbf{递增序列}（或不降序列），并定义
\[
\lim_{n\to\infty}A_n:=\bigcup_{n=1}^{\infty}A_n.
\]
若
\[
A_1\supseteq A_2\supseteq\cdots\supseteq A_n\supseteq\cdots,
\]
则称其为\textbf{递减序列}（或不增序列），并定义
\[
\lim_{n\to\infty}A_n:=\bigcap_{n=1}^{\infty}A_n.
\]
\end{mydefinition}
\origin{原笔记第1.2节，定义“事件序列的极限”；源码第720行。}

\begin{mytheorem}{概率的连续性}{continuity-of-probability}
\begin{enumerate}
\item 若 $\{A_n\}_{n\geq1}$ 为递增序列，则
\[
P\left(\lim_{n\to\infty}A_n\right)=\lim_{n\to\infty}P(A_n).
\]
称为概率的\textbf{下连续性}；
\item 若 $\{A_n\}_{n\geq1}$ 为递减序列，则
\[
P\left(\lim_{n\to\infty}A_n\right)=\lim_{n\to\infty}P(A_n).
\]
称为概率的\textbf{上连续性}。
\end{enumerate}
\end{mytheorem}
\origin{原笔记第1.2节，定理“概率的连续性”；源码第739行。}

\begin{proposition}{可列次可加性（布尔不等式）}{countable-subadditivity}
对任意可列个事件 $A_1,A_2,\ldots$，有
\[
P\left(\bigcup_{i=1}^{\infty}A_i\right)
\leq\sum_{i=1}^{\infty}P(A_i).
\]
\end{proposition}
\origin{原笔记第1.2节，命题“可列次可加性（布尔不等式）”；源码第779行。}

\subsection{例题：可列个零概率事件的并}\label{note:c1:e6}
\paragraph{题目}
证明：若事件列 $\{A_i\}_{i\geq1}$ 满足 $P(A_i)=0$（$i\geq1$），则
\[
P\left(\bigcup_{i=1}^{\infty}A_i\right)=0.
\]

\basis{原笔记第1.2节，例“可列个零概率事件的并”；源码\texttt{chapters/chapter1.tex}第803行。}
\begin{proof}
由布尔不等式，
\[
0\leq P\left(\bigcup_{i=1}^{\infty}A_i\right)
\leq\sum_{i=1}^{\infty}P(A_i)=0.
\]
\end{proof}

\begin{mydefinition}{事件列的上极限与下极限}{limsup-liminf-events}
设 $A_n\in\mathscr{F}$（$n\geq1$），不要求事件列 $\{A_n\}$ 单调。定义
\[
\limsup_{n\to\infty}A_n
:=\bigcap_{n=1}^{\infty}\bigcup_{k=n}^{\infty}A_k
\]
为事件列的\textbf{上极限}，定义
\[
\liminf_{n\to\infty}A_n
:=\bigcup_{n=1}^{\infty}\bigcap_{k=n}^{\infty}A_k
\]
为事件列的\textbf{下极限}。
\end{mydefinition}
\origin{原笔记第1.2节，定义“事件列的上极限与下极限”；源码第824行。}

\begin{mydefinition}{事件列极限的存在}{event-limit-exists}
	当一个事件列的上极限与下极限相等时，也即
	\[\liminf_{n\to\infty}A_n =\limsup_{n\to\infty}A_n\]
	时，我们定义这个事件列的极限为
	\[\lim\limits_{n\to\infty}A_n :=\liminf_{n\to\infty}A_n =\limsup_{n\to\infty}A_n\]
\end{mydefinition}
\origin{原笔记第1.2节，定义“事件列极限的存在”；源码第860行。}

\subsection{例题：伯努利概率空间}\label{note:c1:e7}
\paragraph{题目}
设 $\Omega$ 是非空集合，$\varnothing\subsetneq A\subsetneq\Omega$，并令
\[
\mathscr{F}=\{\varnothing,\Omega,A,A^c\}.
\]
给定 $0<p<1$，定义
\[
P(\varnothing)=0,\qquad P(\Omega)=1,
\qquad P(A)=p,\qquad P(A^c)=1-p.
\]
验证 $(\Omega,\mathscr F,P)$ 是概率空间，并说明此模型表示的试验。

\basis{原笔记第1.2节，例“伯努利概率空间”；源码\texttt{chapters/chapter1.tex}第868行。}
\begin{proof}[验证]
$\mathscr F$ 是由 $A,A^c$ 的所有并集构成的 $\sigma$-域。所给四个概率均非负，且 $P(\Omega)=1$。

任取 $\mathscr F$ 中两两互斥的可列事件，非空事件至多有两个。若只有一个非空事件，可列可加性直接成立；若有两个，则只能为 $A$ 与 $A^c$，且
\[P(A\cup A^c)=1=p+(1-p)=P(A)+P(A^c).\]
若全部为空集，两边均为零。因此 $P$ 满足可列可加性，$(\Omega,\mathscr F,P)$ 构成概率空间，称为\textbf{伯努利概率空间}（或伯努利概率模型），用于讨论单个事件 $A$ 是否发生。
\end{proof}

\subsection{例题：古典概率空间}\label{note:c1:e8}
\paragraph{题目}
设 $\Omega$ 是只含有限个样本点的非空集合，并取
\[
\mathscr{F}=2^\Omega=\{A:A\subseteq\Omega\}.
\]
若 $n(A)$ 表示集合 $A$ 中元素的个数，定义
\[
P(A)=\frac{n(A)}{n(\Omega)},\qquad A\in\mathscr{F},
\]
验证 $(\Omega,\mathscr F,P)$ 是概率空间，并求每个单点事件的概率。

\basis{原笔记第1.2节，例“古典概率空间”；源码\texttt{chapters/chapter1.tex}第881行。}
\begin{proof}[验证]
$\mathscr F=2^\Omega$ 是 $\Omega$ 上的 $\sigma$-域。由计数的非负性，$P(A)\geq0$；又有 $P(\Omega)=n(\Omega)/n(\Omega)=1$。

若 $A_1,A_2,\ldots$ 两两互斥，由 $\Omega$ 有限，至多有限个 $A_i$ 非空。各非空集合的元素不重复，故
\[P\left(\bigcup_{i=1}^\infty A_i\right)=\frac{n(\bigcup_i A_i)}{n(\Omega)}=\sum_{i=1}^\infty\frac{n(A_i)}{n(\Omega)}=\sum_{i=1}^\infty P(A_i).\]
因此它构成古典概率空间。每个 $\omega\in\Omega$ 都满足
\[P(\{\omega\})=\frac1{n(\Omega)}.\]
\end{proof}

\subsection{例题：最大号码问题}\label{note:c1:e9}
\paragraph{题目}
某城市有编号为 $1,2,\ldots,N$ 的 $N$ 部卡车。独立记录遇到的 $n$ 辆卡车的编号，允许重复，且每次记录到各编号的概率均为 $1/N$。求记录到的最大编号为 $k$（$1\leq k\leq N$）的概率。

\basis{原笔记第1.2节，例“最大号码问题”；源码\texttt{chapters/chapter1.tex}第895行。}
\begin{proof}[解答]
令 $B_k$ 表示最大编号恰为 $k$，$A_k$ 表示最大编号不超过 $k$。则
\[
B_1=A_1,\qquad B_k=A_k\setminus A_{k-1}\quad(2\leq k\leq N),
\]
且
\[
P(A_k)=\left(\frac{k}{N}\right)^n.
\]
因此
\[
P(B_1)=\frac{1}{N^n},
\]
而当 $2\leq k\leq N$ 时，
\[
P(B_k)=P(A_k)-P(A_{k-1})
=\frac{k^n-(k-1)^n}{N^n}.
\]
\end{proof}

\subsection{例题：配对问题}\label{note:c1:e10}
\paragraph{题目}
在有 $n$ 人参加的晚会上，每人带来一顶各不相同的帽子。离场时，每人从随机混放的帽子中无放回地抽取一顶，所有分配方式等可能。求：
\begin{enumerate}
\item 至少有一人抽到自己帽子的概率。
\item 当 $n\to\infty$ 时，至少有一人抽到自己帽子与无人抽到自己帽子的概率极限。
\end{enumerate}

\basis{原笔记第1.2节，例“配对问题”；源码\texttt{chapters/chapter1.tex}第920行。}
\begin{proof}[解答]
\begin{enumerate}
\item 给人和帽子分别编号 $1,2,\ldots,n$，令 $A_i$ 表示第 $i$ 人抽到自己的帽子。对任意
$1\leq i_1<i_2<\cdots<i_r\leq n$，有
\[
P(A_{i_1}\cap\cdots\cap A_{i_r})=\frac{(n-r)!}{n!}.
\]
同阶交事件共有 $\binom{n}{r}$ 个，故
\[
\sum_{1\leq i_1<\cdots<i_r\leq n}
P(A_{i_1}\cap\cdots\cap A_{i_r})
=\binom{n}{r}\frac{(n-r)!}{n!}=\frac{1}{r!}.
\]
由容斥恒等式，所求概率为
\[
P\left(\bigcup_{i=1}^{n}A_i\right)
=\sum_{r=1}^{n}\frac{(-1)^{r+1}}{r!}
=1-\sum_{r=0}^{n}\frac{(-1)^r}{r!}.
\]
\item 因此
\[
\lim_{n\to\infty}P\left(\bigcup_{i=1}^{n}A_i\right)
=1-e^{-1}\approx0.6321,
\]
而无人抽到自己帽子的概率趋于 $e^{-1}\approx0.3679$。
\end{enumerate}
\end{proof}

\subsection{例题：圆桌夫妻问题}\label{note:c1:e11}
\paragraph{题目}
十对夫妻随机地沿圆桌坐下，求没有妻子坐在丈夫身边的概率。

\basis{原笔记第1.2节，例“圆桌夫妻问题”；源码\texttt{chapters/chapter1.tex}第950行。}
\begin{proof}[解答]
令 $A_i$ 表示第 $i$ 对夫妻相邻（$i=1,2,\ldots,10$）。总坐法数为
$(20-1)!=19!$。指定 $r$ 对夫妻相邻时，将每对夫妻看作一个整体，并计入夫妻内部的两种次序，故
\[
P(A_{i_1}\cap\cdots\cap A_{i_r})
=\frac{2^r(20-r-1)!}{19!}
=\frac{2^r(19-r)!}{19!}.
\]
由容斥恒等式，至少有一对夫妻相邻的概率为
\[
P\left(\bigcup_{i=1}^{10}A_i\right)
=\sum_{r=1}^{10}(-1)^{r+1}\binom{10}{r}
\frac{2^r(19-r)!}{19!}
\approx0.6605.
\]
所以，没有夫妻相邻的概率约为
\[
1-0.6605=0.3395.
\]
\end{proof}

\subsection{例题：反复抽取扑克牌中的大王}\label{note:c1:e12}
\paragraph{题目}
从一副 $54$ 张的扑克牌中反复抽牌，每次等可能地抽取一张，记录后放回并混匀，各次抽取相互独立。求：
\begin{enumerate}
\item 一直抽不到大王的概率；
\item 无穷多次抽到大王的概率。
\end{enumerate}

\basis{原笔记第1.2节，例“反复抽取扑克牌中的大王”；源码\texttt{chapters/chapter1.tex}第975行。}
\begin{proof}[解答]
设 $E_n$ 表示第 $n$ 次抽到大王。

\begin{enumerate}
\item 令
\[
C_n=\bigcap_{k=1}^{n}E_k^c.
\]
则 $C_n$ 表示前 $n$ 次均未抽到大王，且 $\{C_n\}$ 为递减序列。令
\[
C=\bigcap_{n=1}^{\infty}C_n,
\]
则 $C$ 表示始终抽不到大王。由独立性和概率的上连续性，
\[
P(C)=\lim_{n\to\infty}P(C_n)
=\lim_{n\to\infty}\left(\frac{53}{54}\right)^n=0.
\]
因此，最终至少抽到一次大王的概率为 $1$。
\item 令 $B$ 表示无穷多次抽到大王，则
\[
B=\limsup_{n\to\infty}E_n
=\bigcap_{n=1}^{\infty}\bigcup_{k=n}^{\infty}E_k.
\]
再令
\[
B_n=\bigcup_{k=n}^{\infty}E_k.
\]
事件 $B_n$ 表示“从第 $n$ 次抽牌起至少抽到一次大王”。由第一个问题同样的计算，
$P(B_n^c)=0$，故 $P(B_n)=1$。由于 $\{B_n\}$ 为递减序列，
\[
P(B)=P\left(\bigcap_{n=1}^{\infty}B_n\right)
=\lim_{n\to\infty}P(B_n)=1.
\]
\end{enumerate}
\end{proof}

\begin{mydefinition}{测度、有限测度与 $\sigma$ 有限测度}{measure-space}
设 $\Omega$ 是非空集合，$\mathscr{F}$ 是 $\Omega$ 上给定的一个 $\sigma$-域。若映射
\[
\mu:\mathscr{F}\longrightarrow[0,+\infty]
\]
满足：
\begin{enumerate}
\item $\mu(\varnothing)=0$；
\item 对任意两两互斥的事件列 $A_1,A_2,\ldots\in\mathscr{F}$，有
\[
\mu\left(\bigcup_{n=1}^{\infty}A_n\right)
=\sum_{n=1}^{\infty}\mu(A_n),
\]
\end{enumerate}
则称 $\mu$ 为 $\mathscr{F}$ 上的\textbf{测度}，三元组 $(\Omega,\mathscr{F},\mu)$ 称为\textbf{测度空间}。

\end{mydefinition}
\origin{原笔记第1.2节，定义“测度、有限测度与 $\sigma$ 有限测度”；源码第1024行。}

若存在两两互斥的事件列 $A_1,A_2,\ldots\in\mathscr{F}$，使得
\[
\bigcup_{n=1}^{\infty}A_n=\Omega,
\qquad
\mu(A_n)<+\infty\quad(n\geq1),
\]
则称 $\mu$ 为\textbf{$\sigma$ 有限测度}。若 $\mu(\Omega)<+\infty$，则称 $\mu$ 为\textbf{有限测度}；若 $\mu(\Omega)=1$，则称 $\mu$ 为\textbf{概率测度}，也简称为\textbf{概率}。

\begin{proposition}{测度的有限可加性}{measure-finite-additivity}
若 $A_1,\ldots,A_n\in\mathscr{F}$ 两两互斥，则
\[
\mu\left(\bigcup_{i=1}^{n}A_i\right)=\sum_{i=1}^{n}\mu(A_i).
\]
\end{proposition}
\origin{原笔记第1.2节，命题“测度的有限可加性”；源码第1050行。}

\begin{proposition}{测度的单调性}{measure-monotonicity}
若 $A,B\in\mathscr{F}$ 且 $B\subseteq A$，则
\[
\mu(B)\leq\mu(A).
\]
\end{proposition}
\origin{原笔记第1.2节，命题“测度的单调性”；源码第1061行。}

\begin{proposition}{测度的可减性}{measure-subtractivity}
若 $A,B\in\mathscr{F}$，$A\subseteq B$，且 $\mu(A)<+\infty$，则
\[
\mu(B\setminus A)=\mu(B)-\mu(A).
\]
\end{proposition}
\origin{原笔记第1.2节，命题“测度的可减性”；源码第1075行。}

\begin{proposition}{测度的容斥恒等式}{measure-inclusion-exclusion}
若 $A_1,A_2,\ldots,A_n\in\mathscr{F}$，且对每个 $1\leq k\leq n$ 都有 $\mu(A_k)<+\infty$，则
\[
\mu\left(\bigcup_{k=1}^{n}A_k\right)
=\sum_{r=1}^{n}(-1)^{r+1}
\sum_{1\leq i_1<i_2<\cdots<i_r\leq n}
\mu(A_{i_1}\cap A_{i_2}\cap\cdots\cap A_{i_r}).
\]
\end{proposition}
\origin{原笔记第1.2节，命题“测度的容斥恒等式”；源码第1090行。}

\begin{mytheorem}{测度的可列次可加性}{measure-countable-subadditivity}
对任意事件列 $A_1,A_2,\ldots\in\mathscr{F}$，有
\[
\mu\left(\bigcup_{n=1}^{\infty}A_n\right)
\leq\sum_{n=1}^{\infty}\mu(A_n).
\]
\end{mytheorem}
\origin{原笔记第1.2节，定理“测度的可列次可加性”；源码第1104行。}

\begin{proposition}{测度的下连续性}{measure-lower-continuity}
若 $\{A_n\}_{n\geq1}\subseteq\mathscr{F}$ 为递增序列，则
\[
\mu\left(\lim_{n\to\infty}A_n\right)
=\lim_{n\to\infty}\mu(A_n).
\]
\end{proposition}
\origin{原笔记第1.2节，命题“测度的下连续性”；源码第1129行。}

\begin{proposition}{测度的上连续性}{measure-upper-continuity}
若 $\{A_n\}_{n\geq1}\subseteq\mathscr{F}$ 为递减序列，且 $\mu(A_1)<+\infty$，则
\[
\mu\left(\lim_{n\to\infty}A_n\right)
=\lim_{n\to\infty}\mu(A_n).
\]
\end{proposition}
\origin{原笔记第1.2节，命题“测度的上连续性”；源码第1141行。}

\subsection{例题：Lebesgue 测度}\label{note:c1:e13}
\paragraph{题目}
在 $\mathscr{B}(\mathbb{R}^d)$ 上存在唯一的测度 $\lambda$，使得对任意半开长方体
\[
(a_1,b_1]\times(a_2,b_2]\times\cdots\times(a_d,b_d]
\]
（其中 $a_i\leq b_i$，$1\leq i\leq d$），都有
\[
\lambda\left((a_1,b_1]\times\cdots\times(a_d,b_d]\right)
=(b_1-a_1)(b_2-a_2)\cdots(b_d-a_d).
\]
该测度称为 $\mathbb{R}^d$ 上的\textbf{Lebesgue 测度}，它给出了通常所说的长度、面积和体积。以该测度及其长方体公式为已知，验证它在 $d=1,2,3$ 时分别给出半开区间、矩形和长方体的通常长度、面积和体积。

\basis{原笔记第1.2节，例“Lebesgue 测度”；源码\texttt{chapters/chapter1.tex}第1167行。}
\begin{proof}[验证]
由题设公式，当 $d=1$ 时，
\[\lambda((a_1,b_1])=b_1-a_1,\]
正是区间的长度。当 $d=2$ 时，
\[\lambda((a_1,b_1]\times(a_2,b_2])=(b_1-a_1)(b_2-a_2),\]
正是矩形面积。当 $d=3$ 时，
\[\lambda((a_1,b_1]\times(a_2,b_2]\times(a_3,b_3])=(b_1-a_1)(b_2-a_2)(b_3-a_3),\]
正是长方体体积。若某一边长为零，公式同样给出零。
\end{proof}

\clearpage\section{条件概率}\label{k:conditional}

\subsection{例题：掷两枚骰子}\label{note:c1:e14}
\paragraph{题目}
独立地依次掷两枚均匀骰子。已知第一枚骰子的点数为 $4$，求两枚骰子点数之和为偶数的概率。

\basis{原笔记第1.3节，例“掷两枚骰子”；源码\texttt{chapters/chapter1.tex}第1186行。}
\begin{proof}[解答]
设 $B$ 表示“第一枚骰子的点数为 $4$”，$A$ 表示“两枚骰子点数之和为偶数”。在 $B$ 已经发生时，第二枚骰子共有 $6$ 种等可能结果；其中点数为 $2,4,6$ 时 $A$ 发生，故
\[
P(A\mid B)=\frac{3}{6}=\frac12.
\]
\end{proof}

\begin{mydefinition}{条件概率}{conditional-probability}
设 $(\Omega,\mathscr{F},P)$ 为概率空间，$A,B\in\mathscr{F}$，且 $P(B)>0$。称
\[
P(A\mid B)=\frac{P(A\cap B)}{P(B)}
\]
为在事件 $B$ 发生的条件下事件 $A$ 的\textbf{条件概率}。
\end{mydefinition}
\origin{原笔记第1.3节，定义“条件概率”；源码第1195行。}

\subsection{例题：寻找钥匙}\label{note:c1:e15}
\paragraph{题目}
一串钥匙可能落在左口袋或右口袋。设 $L$ 表示“钥匙在左口袋”，$R$ 表示“钥匙在右口袋”，且
\[
P(L)=0.4,\qquad P(R)=0.4,\qquad R\subseteq L^{\mathrm c}.
\]
检查左口袋后发现没有钥匙，即已知 $L^{\mathrm c}$ 发生。求此时钥匙在右口袋的条件概率 $P(R\mid L^{\mathrm c})$。

\basis{原笔记第1.3节，例“寻找钥匙”；源码\texttt{chapters/chapter1.tex}第1203行。}
\begin{proof}[解答]
由 $R\subseteq L^{\mathrm c}$，有 $R\cap L^{\mathrm c}=R$，因此
\[
P(R\mid L^{\mathrm c})
=\frac{P(R\cap L^{\mathrm c})}{P(L^{\mathrm c})}
=\frac{P(R)}{1-P(L)}
=\frac{0.4}{0.6}
=\frac23.
\]
\end{proof}

\begin{mytheorem}{条件概率的乘法公式}{conditional-multiplication-formula}
若 $A,B\in\mathscr{F}$ 且 $P(B)>0$，则
\[
P(A\cap B)=P(B)P(A\mid B).
\]

更一般地，若 $A_1,\ldots,A_n\in\mathscr{F}$ 且
\[
P(A_1\cap\cdots\cap A_{n-1})>0,
\]
则
\[
\begin{aligned}
P(A_1\cap\cdots\cap A_n)
={}&P(A_1)P(A_2\mid A_1)P(A_3\mid A_1\cap A_2)\cdots\\
&\cdot P(A_n\mid A_1\cap\cdots\cap A_{n-1}).
\end{aligned}
\]
\end{mytheorem}
\origin{原笔记第1.3节，定理“条件概率的乘法公式”；源码第1220行。}

\begin{mytheorem}{全概率公式}{total-probability-formula}
设 $A_1,A_2,\ldots$ 构成样本空间 $\Omega$ 的一个可列分割，即
\[
A_i\cap A_j=\varnothing\quad(i\ne j),
\qquad
\bigcup_{i=1}^{\infty}A_i=\Omega,
\]
并且 $P(A_i)>0$。则对任意 $B\in\mathscr{F}$，
\[
P(B)=\sum_{i=1}^{\infty}P(A_i)P(B\mid A_i).
\]
当分割只有有限个事件时，相应地使用有限求和。
\end{mytheorem}
\origin{原笔记第1.3节，定理“全概率公式”；源码第1254行。}

\subsection{例题：同卵双胞胎的比例}\label{note:c1:e16}
\paragraph{题目}
设 $A$ 表示“双胞胎为同卵”，$B$ 表示“双胞胎性别相同”。已知同卵双胞胎性别必然相同，异卵双胞胎两人的性别可视为独立且各以 $1/2$ 的概率为男或女，所以
\[
P(B\mid A)=1,
\qquad
P(B\mid A^{\mathrm c})=\frac12.
\]
总体中双胞胎性别相同的比例为 $P(B)=0.64$，求同卵双胞胎所占比例 $P(A)$。

\basis{原笔记第1.3节，例“同卵双胞胎的比例”；源码\texttt{chapters/chapter1.tex}第1283行。}
\begin{proof}[解答]
若总体中双胞胎性别相同的比例为 $P(B)=0.64$，记 $P(A)=x$。由全概率公式，
\[
0.64=x+\frac12(1-x),
\]
解得
\[
x=0.28.
\]
因此同卵双胞胎约占 $28\%$。
\end{proof}

\subsection{例题：Pólya 罐模型}\label{note:c1:e17}
\paragraph{题目}
罐中原有 $a$ 个红球与 $b$ 个蓝球，其中 $a,b$ 为正整数。每次随机取出一个球，放回后再加入 $c$ 个同色球，其中 $c$ 为非负整数。证明：任意第 $k$ 次取到红球的概率都等于
\[
\frac{a}{a+b}.
\]

\basis{原笔记第1.3节，例“Pólya 罐模型”；源码\texttt{chapters/chapter1.tex}第1301行。}
\begin{proof}
对所有允许的初始红、蓝球数同时作关于 $k$ 的归纳。$k=1$ 时，红球占全部球的比例为 $a/(a+b)$，结论成立。假设对任意初始红、蓝球数，结论对第 $k-1$ 次抽取成立。

以第一次抽取的颜色作分割：若第一次取到红球，则此后罐中有 $a+c$ 个红球与 $b$ 个蓝球；若第一次取到蓝球，则有 $a$ 个红球与 $b+c$ 个蓝球。

由归纳假设和全概率公式，
\[
\begin{aligned}
P(\text{第 $k$ 次为红球})
&=\frac{a}{a+b}\cdot\frac{a+c}{a+b+c}
+\frac{b}{a+b}\cdot\frac{a}{a+b+c}\\
&=\frac{a}{a+b}.
\end{aligned}
\]
故结论对所有 $k\geq1$ 成立。
\end{proof}

\subsection{例题：\texorpdfstring{配对问题的递推与恰有 $k$ 个配对}{配对问题的递推与恰有 k 个配对}}\label{note:c1:e18}
\paragraph{题目}
将 $n$ 封信等可能地一一放入写有相应地址的 $n$ 个信封中。令 $E_n$ 表示没有一封信放入正确信封，记
\[q_n=P(E_n),\qquad q_0=1,\quad q_1=0.\]
\begin{enumerate}
\item 证明当 $n\geq2$ 时
\[q_n=\frac1n q_{n-2}+\frac{n-1}{n}q_{n-1},\]
并据此求 $q_n$。
\item 求恰有 $k$ 封信放入正确信封的概率，其中 $0\leq k\leq n$。
\end{enumerate}

\basis{原笔记第1.3节，例“配对问题的递推与恰有 $k$ 个配对”；源码\texttt{chapters/chapter1.tex}第1321行。}
\begin{proof}[解答]
\begin{enumerate}
\item 记 $D_n=n!q_n$ 为 $n$ 个元素的错排数。第一封信只能放入其余 $n-1$ 个信封。固定其中一个信封 $j$，将它对应的信记为第 $j$ 封信。

若第 $j$ 封信放入第一封信的信封，则这两封信互换，余下 $n-2$ 封信须全部错排，有 $D_{n-2}$ 种。

若第 $j$ 封信不放入第一封信的信封，则删去第一封信及信封 $j$，并把第一封信的信封改记为信封 $j$。剩余放法与 $n-1$ 封信的错排一一对应：第 $j$ 封信不在改记后的信封 $j$，其他信也不在各自信封；反过来可恢复原放法。因此这一类有 $D_{n-1}$ 种。

两类互斥且穷尽所有错排，所以
\[D_n=(n-1)(D_{n-1}+D_{n-2}).\]
除以 $n!$ 得
\[q_n=\frac1n q_{n-2}+\frac{n-1}{n}q_{n-1}.\]
移项并用 $q_1-q_0=-1$，得到
\begin{align*}
q_n-q_{n-1}&=-\frac1n(q_{n-1}-q_{n-2})=\frac{(-1)^n}{n!},\\
q_n&=q_0+\sum_{r=1}^n(q_r-q_{r-1})=\sum_{r=0}^n\frac{(-1)^r}{r!}.
\end{align*}
\item 先选出放置正确的 $k$ 封信，有 $\binom nk$ 种选择，再将其余 $n-k$ 封信完全错排，有 $D_{n-k}=(n-k)!q_{n-k}$ 种。总放法有 $n!$ 种，因此
\[P(\text{恰有 $k$ 封信放置正确})=\binom nk\frac{(n-k)!q_{n-k}}{n!}=\frac1{k!}\sum_{r=0}^{n-k}\frac{(-1)^r}{r!}.\]
\end{enumerate}
\end{proof}

\begin{mytheorem}{贝叶斯公式}{bayes-formula}
设 $A_1,A_2,\ldots$ 构成 $\Omega$ 的一个可列分割，且 $P(A_i)>0$。若 $B\in\mathscr{F}$ 且 $P(B)>0$，则
\[
P(A_i\mid B)
=\frac{P(A_i)P(B\mid A_i)}
{\displaystyle\sum_{k=1}^{\infty}P(A_k)P(B\mid A_k)}.
\]
有限分割的情形只需将分母改为有限求和。
\end{mytheorem}
\origin{原笔记第1.3节，定理“贝叶斯公式”；源码第1362行。}

\subsection{例题：保险事故原因}\label{note:c1:e19}
\paragraph{题目}
设 $A$ 表示“事故较容易发生”，$B$ 表示“发生事故”。已知
\[
P(A)=0.3,
\qquad
P(B\mid A)=0.4,
\qquad
P(B\mid A^{\mathrm c})=0.2.
\]
求：
\begin{enumerate}
\item 发生事故的概率 $P(B)$。
\item 已发生事故时，属于事故较容易发生这一情形的条件概率 $P(A\mid B)$。
\end{enumerate}

\basis{原笔记第1.3节，例“保险事故原因”；源码\texttt{chapters/chapter1.tex}第1384行。}
\begin{proof}[解答]
\begin{enumerate}
\item 由全概率公式，
\[
P(B)=0.3\times0.4+0.7\times0.2=0.26.
\]
\item 再由贝叶斯公式，
\[
P(A\mid B)
=\frac{0.3\times0.4}{0.26}
=\frac6{13}
\approx0.4615.
\]
\end{enumerate}
\end{proof}

\subsection{例题：Monty Hall 问题}\label{note:c1:e20}
\paragraph{题目}
汽车以相同概率位于三扇门中的任意一扇门后；三扇门后分别藏有一辆汽车和两只山羊。参赛者先选择第 $1$ 扇门；主持人知道汽车的位置，随后一定在第 $2,3$ 扇门中打开一扇藏有山羊的门。若两扇门都可打开，主持人以相同概率选择其中一扇。已知主持人打开了第 $3$ 扇门，求不换门与换门的获胜概率，并判断是否应改选第 $2$ 扇门。

\basis{原笔记第1.3节，例“Monty Hall 问题”；源码\texttt{chapters/chapter1.tex}第1406行。}
\begin{proof}[解答]
令 $A_i$ 表示“汽车在第 $i$ 扇门后”，$B$ 表示“主持人打开第 $3$ 扇门”。则
\[
P(A_1)=P(A_2)=P(A_3)=\frac13,
\]
且
\[
P(B\mid A_1)=\frac12,
\qquad
P(B\mid A_2)=1,
\qquad
P(B\mid A_3)=0.
\]
因此
\[
P(B)=\frac13\cdot\frac12+\frac13\cdot1+\frac13\cdot0=\frac12.
\]
由贝叶斯公式，
\[
P(A_1\mid B)=\frac13,
\qquad
P(A_2\mid B)=\frac23.
\]
所以改选第 $2$ 扇门的获胜概率为 $2/3$，应当换门。
\end{proof}

\begin{mytheorem}{条件概率构成概率测度}{conditional-probability-measure}
固定 $B\in\mathscr{F}$ 且 $P(B)>0$，定义
\[
P_B(A)=P(A\mid B),\qquad A\in\mathscr{F}.
\]
则 $P_B$ 是可测空间 $(\Omega,\mathscr{F})$ 上的一个概率测度。
\end{mytheorem}
\origin{原笔记第1.3节，定理“条件概率构成概率测度”；源码第1436行。}

\begin{mydefinition}{条件概率空间中的再次条件化}{iterated-conditioning}
固定 $B\in\mathscr{F}$ 且 $P(B)>0$。若 $C\in\mathscr{F}$ 且 $P_B(C)>0$，定义
\[
P_B(A\mid C)=\frac{P_B(A\cap C)}{P_B(C)}.
\]
\end{mydefinition}
\origin{原笔记第1.3节，定义“条件概率空间中的再次条件化”；源码第1463行。}

\begin{proposition}{再次条件化等价于对交事件条件化}{nested-conditioning}
固定 $B\in\mathscr{F}$ 且 $P(B)>0$。对任意 $C\in\mathscr{F}$，有
\[
P_B(C)>0\quad\Longleftrightarrow\quad P(B\cap C)>0.
\]
在上述条件成立时，对任意 $A\in\mathscr{F}$，
\[
P_B(A\mid C)=P(A\mid B\cap C).
\]
\end{proposition}
\origin{原笔记第1.3节，命题“再次条件化等价于对交事件条件化”；源码第1470行。}

\subsection{例题：血液化验与重复检测}\label{note:c1:e21}
\paragraph{题目}
某项血液化验对患者有 $95\%$ 的概率检出阳性，对健康人有 $1\%$ 的概率出现伪阳性，患者占总人口的 $0.5\%$。记 $A$ 为患病，$B$ 为第一次阳性，$B_2$ 为第二次阳性。
\begin{enumerate}
\item 若一次化验呈阳性，求此人确实患病的条件概率 $P(A\mid B)$。
\item 再次接受同样化验，假设两次化验在给定患病状态后条件独立，即
\[P(B_2\mid A\cap B)=P(B_2\mid A)=0.95,\qquad P(B_2\mid A^{\mathrm c}\cap B)=P(B_2\mid A^{\mathrm c})=0.01.\]
求连续两次呈阳性后的患病概率 $P(A\mid B\cap B_2)$。
\end{enumerate}

\basis{原笔记第1.3节，例“血液化验与重复检测”；源码\texttt{chapters/chapter1.tex}第1515行。}
\begin{proof}[解答]
\begin{enumerate}
\item 设 $A$ 表示“此人患病”，$B$ 表示“第一次化验呈阳性”。则
\[
P(A)=0.005,
\qquad
P(B\mid A)=0.95,
\qquad
P(B\mid A^{\mathrm c})=0.01.
\]
由贝叶斯公式，
\[
\begin{aligned}
P(A\mid B)
&=\frac{P(A)P(B\mid A)}
{P(A)P(B\mid A)+P(A^{\mathrm c})P(B\mid A^{\mathrm c})}\\
&=\frac{0.005\times0.95}
{0.005\times0.95+(1-0.005)\times0.01}\\
&\approx0.323.
\end{aligned}
\]
因此，即使检验对患者的阳性检出率很高，一次阳性结果所对应的患病概率也只有约 $32.3\%$。原因在于患病率很低，健康人数量很大，伪阳性在全部阳性结果中占据了不可忽略的比例。
\item 由题设的条件独立性，连续两次呈阳性后，
\[
\begin{aligned}
P(A\mid B\cap B_2)
&=\frac{P(A)P(B\mid A)P(B_2\mid A)}
{P(A)P(B\mid A)P(B_2\mid A)
+P(A^{\mathrm c})P(B\mid A^{\mathrm c})P(B_2\mid A^{\mathrm c})}\\
&=\frac{0.005\times0.95^2}
{0.005\times0.95^2+0.995\times0.01^2}\\
&\approx0.978.
\end{aligned}
\]
连续两次阳性使后验患病概率提高到约 $97.8\%$。
\end{enumerate}
\end{proof}

\subsection{例题：飞机失事搜索}\label{note:c1:e22}
\paragraph{题目}
一架飞机失踪，并且等可能地落在第 $1,2,3$ 号区域。设 $B_i$ 表示“飞机落在第 $i$ 个区域”，$A_i$ 表示“搜索第 $i$ 个区域后没有找到飞机”，并设
\[
P(A_i\mid B_i)=\beta_i,
\qquad 0<\beta_i<1,
\qquad i=1,2,3.
\]
若飞机不在被搜索的区域，则在那里搜索必然找不到飞机。

\begin{enumerate}
\item 只搜索第 $1$ 号区域且未找到飞机，求 $P(B_i\mid A_1)$（$i=1,2,3$）。
\item 随后搜索第 $2$ 号区域仍未找到飞机，求 $P(B_i\mid A_1\cap A_2)$（$i=1,2,3$）。
\item 最后搜索第 $3$ 号区域仍未找到飞机，求 $P(B_i\mid A_1\cap A_2\cap A_3)$（$i=1,2,3$）。
\end{enumerate}

\basis{原笔记第1.3节，例“飞机失事搜索”；源码\texttt{chapters/chapter1.tex}第1560行。}
\begin{proof}[解答]
$B_1,B_2,B_3$ 构成完备事件组，且 $P(B_i)=1/3$。若给定 $B_i$，在其他区域搜索必然失败。因此计算多个搜索失败的交事件时，只留下在真实所在区域漏检的条件；不需要假设各次搜索相互独立。
\begin{enumerate}
\item 由 $P(A_1\mid B_1)=\beta_1$、$P(A_1\mid B_2)=P(A_1\mid B_3)=1$，全概率公式给出
\[P(A_1)=\frac{\beta_1+2}{3}.\]
用贝叶斯公式分别计算：
\begin{align*}
P(B_1\mid A_1)&=\frac{(1/3)\beta_1}{(\beta_1+2)/3}=\frac{\beta_1}{\beta_1+2},\\
P(B_2\mid A_1)&=P(B_3\mid A_1)=\frac{(1/3)\cdot1}{(\beta_1+2)/3}=\frac1{\beta_1+2}.
\end{align*}
\item 给定 $B_1$，事件 $A_2$ 必然发生，故 $P(A_1\cap A_2\mid B_1)=\beta_1$；给定 $B_2$，事件 $A_1$ 必然发生，故该条件概率为 $\beta_2$；给定 $B_3$，两次搜索均必然失败，故该条件概率为 $1$。于是
\[P(A_1\cap A_2)=\frac{\beta_1+\beta_2+1}{3}.\]
由贝叶斯公式，
\begin{align*}
P(B_1\mid A_1\cap A_2)&=\frac{(1/3)\beta_1}{(\beta_1+\beta_2+1)/3}=\frac{\beta_1}{\beta_1+\beta_2+1},\\
P(B_2\mid A_1\cap A_2)&=\frac{\beta_2}{\beta_1+\beta_2+1},\\
P(B_3\mid A_1\cap A_2)&=\frac1{\beta_1+\beta_2+1}.
\end{align*}
\item 给定 $B_i$，另外两个区域的搜索失败是必然事件，故
\[P(A_1\cap A_2\cap A_3\mid B_i)=P(A_i\mid B_i)=\beta_i.\]
因此
\[P(A_1\cap A_2\cap A_3)=\frac{\beta_1+\beta_2+\beta_3}{3},\]
进而
\[P(B_i\mid A_1\cap A_2\cap A_3)=\frac{(1/3)\beta_i}{(\beta_1+\beta_2+\beta_3)/3}=\frac{\beta_i}{\beta_1+\beta_2+\beta_3},\quad i=1,2,3.\]
每次搜索失败后，都可将当前后验概率作为下一次更新的先验概率。
\end{enumerate}
\end{proof}

\subsection{例题：赌徒破产模型}\label{note:c1:e23}
\paragraph{题目}
设 $a$ 为正整数，质点从整数点 $i\in\{0,1,\ldots,a\}$ 出发。只要质点位于 $1,2,\ldots,a-1$ 中，它就以概率 $p$ 向右移动一步，以概率 $q=1-p$ 向左移动一步，其中 $0<p<1$；一旦到达 $0$ 或 $a$，便停留在那里。

\begin{enumerate}
\item 分别求质点最终到达 $0$ 与到达 $a$ 的概率，并说明它是否以概率 $1$ 被端点吸收。
\item 将位置解释为甲的赌本：甲、乙初始赌本为 $i$ 元和 $a-i$ 元，每局胜者从败者处获得 $1$ 元，甲每局获胜概率为 $p$。固定 $i$，令 $a\to\infty$，求甲最终破产概率的极限。
\end{enumerate}

\basis{原笔记第1.3节，例“赌徒破产模型”；源码\texttt{chapters/chapter1.tex}第1594行。}
\begin{proof}[解答]
\begin{enumerate}
\item 设 $E_i$ 表示“质点从 $i$ 出发”，$F$ 表示“最终到达 $0$”，并记
\[
Q_i=P(F\mid E_i).
\]
按第一次移动的方向使用全概率公式，得到
\[
Q_i=pQ_{i+1}+qQ_{i-1},
\qquad 1\leq i\leq a-1,
\]
边界条件为
\[
Q_0=1,
\qquad
Q_a=0.
\]
递推式可改写为
\[
Q_{i+1}-Q_i=\frac{q}{p}(Q_i-Q_{i-1}).
\]

令 $d_j=Q_j-Q_{j-1}$。当 $p=q=1/2$ 时，$d_j=d_1$，故
\[Q_i=1+id_1,\qquad 0=Q_a=1+ad_1.\]
由此得到
\[
Q_i=\frac{a-i}{a}.
\]
当 $p\ne q$ 时，令 $\rho=q/p\ne1$，则 $d_j=\rho^{j-1}d_1$。对差分求和，
\begin{align*}
Q_i&=1+d_1\sum_{j=1}^i\rho^{j-1}=1+d_1\frac{1-\rho^i}{1-\rho},\\
0=Q_a&=1+d_1\frac{1-\rho^a}{1-\rho}.
\end{align*}
消去 $d_1$，得到
\[
Q_i
=\frac{\left(\dfrac qp\right)^i-\left(\dfrac qp\right)^a}
{1-\left(\dfrac qp\right)^a}.
\]
因此
\[
P(F\mid E_i)=
\begin{cases}
\dfrac{\left(\dfrac qp\right)^i-\left(\dfrac qp\right)^a}
{1-\left(\dfrac qp\right)^a},&p\ne q,\\[1.2em]
\dfrac{a-i}{a},&p=q=\dfrac12.
\end{cases}
\]

同理，若 $M$ 表示“最终到达 $a$”，则
\[
P(M\mid E_i)=
\begin{cases}
\dfrac{1-\left(\dfrac qp\right)^i}
{1-\left(\dfrac qp\right)^a},&p\ne q,\\[1.2em]
\dfrac{i}{a},&p=q=\dfrac12.
\end{cases}
\]
两式相加得到
\[
P(F\cup M\mid E_i)=P(F\mid E_i)+P(M\mid E_i)=1,
\]
即质点以概率 $1$ 被某一个端点吸收。
\item 这个模型也可解释为赌徒破产问题：甲、乙的初始赌本分别为 $i$ 元和 $a-i$ 元，每局胜者从败者处获得 $1$ 元，甲每局获胜的概率为 $p$。固定甲的初始赌本 $i$，令乙的初始赌本趋于无穷，即令 $a\to\infty$，则
\[
\lim_{a\to\infty}P(F\mid E_i)=
\begin{cases}
\left(\dfrac qp\right)^i,&p>\dfrac12,\\[0.8em]
1,&p\leq\dfrac12.
\end{cases}
\]
因此，当 $i>0$ 且 $p>1/2$ 时，甲有正概率永不破产；当 $p\leq1/2$ 时，面对拥有无限赌资的对手，甲以概率 $1$ 破产。
\end{enumerate}
\end{proof}

\clearpage\section{事件的独立性和试验的独立性}\label{k:independence}

\begin{mydefinition}{两个事件的独立性}{independence-two-events}
设 $(\Omega,\mathscr{F},P)$ 为概率空间，$A,B\in\mathscr{F}$。若
\[
P(A\cap B)=P(A)P(B),
\]
则称事件 $A$ 与 $B$\textbf{相互独立}。
\end{mydefinition}
\origin{原笔记第1.4节，定义“两个事件的独立性”；源码第1675行。}

\begin{proposition}{独立性对补事件保持}{independence-complements}
若事件 $A$ 与 $B$ 相互独立，则下面三对事件也分别相互独立：
\[
A\text{ 与 }B^{\mathrm c},
\qquad
A^{\mathrm c}\text{ 与 }B,
\qquad
A^{\mathrm c}\text{ 与 }B^{\mathrm c}.
\]
\end{proposition}
\origin{原笔记第1.4节，命题“独立性对补事件保持”；源码第1691行。}

\subsection{例题：独立与互斥的区别}\label{note:c1:e24}
\paragraph{题目}
向区间 $[0,1]$ 内均匀随机投入一个粒子，以下区间表示粒子落入其中的事件。分别判断每对事件是否互斥、是否相互独立：
\begin{enumerate}
\item \[
A_1=\left[0,\frac12\right],
\qquad
B_1=\left[\frac14,\frac34\right].
\]
\item \[
A_2=\left[\frac15,\frac13\right],
\qquad
B_2=\left[\frac14,1\right].
\]
\end{enumerate}

\basis{原笔记第1.4节，例“独立与互斥的区别”；源码\texttt{chapters/chapter1.tex}第1726行。}
\begin{proof}[解答]
\begin{enumerate}
\item 有
\[
A_1\cap B_1=\left[\frac14,\frac12\right],
\]
并且
\[
P(A_1)=\frac12,
\qquad
P(B_1)=\frac12,
\qquad
P(A_1\cap B_1)=\frac14=P(A_1)P(B_1).
\]
所以 $A_1$ 与 $B_1$ 不互斥，但相互独立。
\item 有
\[
A_2\cap B_2=\left[\frac14,\frac13\right],
\]
且
\[
P(A_2)=\frac{2}{15},
\qquad
P(B_2)=\frac34,
\qquad
P(A_2\cap B_2)=\frac1{12}\ne\frac1{10}=P(A_2)P(B_2).
\]
所以 $A_2$ 与 $B_2$ 既不互斥，也不相互独立。
\end{enumerate}
\end{proof}

\begin{mydefinition}{两两独立与相互独立}{pairwise-mutual-independence}
	设 $A_1,A_2,\ldots,A_n$ 为 $n\geq2$ 个事件。

	若对任意 $i\ne j$ 都有
	\[
	P(A_i\cap A_j)=P(A_i)P(A_j),
	\]
	则称 $A_1,A_2,\ldots,A_n$\textbf{两两独立}。

	若对任意 $k\in\{2,3,\ldots,n\}$ 以及任意不同的指标
	\[
	1\leq i_1<i_2<\cdots<i_k\leq n,
	\]
	都有
	\[
	P(A_{i_1}\cap A_{i_2}\cap\cdots\cap A_{i_k})
	=\prod_{r=1}^{k}P(A_{i_r}),
	\]
	则称 $A_1,A_2,\ldots,A_n$\textbf{相互独立}。
\end{mydefinition}
\origin{原笔记第1.4节，定义“两两独立与相互独立”；源码第1769行。}

\subsection{例题：两两独立不推出相互独立}\label{note:c1:e25}
\paragraph{题目}
独立地掷两枚均匀骰子。令 $E$ 表示“两枚骰子的点数之和为 $7$”，$F$ 表示“第一枚骰子的点数为 $4$”，$G$ 表示“第二枚骰子的点数为 $3$”。判断 $E,F,G$ 是否两两独立、是否相互独立。

\basis{原笔记第1.4节，例“两两独立不推出相互独立”；源码\texttt{chapters/chapter1.tex}第1791行。}
\begin{proof}[解答]
有
\[
P(E)=P(F)=P(G)=\frac16,
\]
并且
\[
P(E\cap F)=P(E\cap G)=P(F\cap G)=P(E\cap F\cap G)=\frac1{36}.
\]
因此
\[
P(E\cap F)=P(E)P(F),
\qquad
P(E\cap G)=P(E)P(G),
\qquad
P(F\cap G)=P(F)P(G),
\]
所以 $E,F,G$ 两两独立。但是
\[
P(E\cap F\cap G)=\frac1{36}
\ne\frac1{216}=P(E)P(F)P(G),
\]
故它们并不相互独立。
\end{proof}

\subsection{例题：三色球模型}\label{note:c1:e26}
\paragraph{题目}
袋中有 $4$ 个球：一个只涂白色，一个只涂蓝色，一个只涂红色，另一个同时涂有白、蓝、红三色。等可能地取出一个球。设 $A,B,C$ 分别表示“取出的球涂有白色、蓝色、红色”。判断 $A,B,C$ 是否两两独立、是否相互独立。

\basis{原笔记第1.4节，例“三色球模型”；源码\texttt{chapters/chapter1.tex}第1816行。}
\begin{proof}[解答]
有
\[
P(A)=P(B)=P(C)=\frac12,
\]
且
\[
P(A\cap B)=P(A\cap C)=P(B\cap C)=P(A\cap B\cap C)=\frac14.
\]
因此
\[
P(A\cap B)=P(A)P(B),
\quad
P(A\cap C)=P(A)P(C),
\quad
P(B\cap C)=P(B)P(C),
\]
所以 $A,B,C$ 两两独立；但
\[
P(A\cap B\cap C)=\frac14
\ne\frac18=P(A)P(B)P(C),
\]
故它们不相互独立。
\end{proof}

\begin{mydefinition}{可列个事件的相互独立}{countable-event-independence}
设 $A_1,A_2,\ldots$ 为可列个事件。若其中任意有限多个不同事件都相互独立，则称事件列 $\{A_n\}_{n\geq1}$\textbf{相互独立}。
\end{mydefinition}
\origin{原笔记第1.4节，定义“可列个事件的相互独立”；源码第1847行。}

\begin{mydefinition}{独立试验、伯努利试验与独立重复试验}{independent-trials}
设一个试验由一系列子试验组成。
\begin{enumerate}
    \item 若分别由不同子试验决定的任意有限个事件相互独立，则称这些子试验\textbf{相互独立}。
    \item 若一个试验由 $n$ 个相互独立且同分布的子试验组成，则称其为\textbf{$n$ 重独立试验}。若每个子试验只有“成功”和“失败”两种结果，且成功概率为 $p$，$0<p<1$，则称其为\textbf{$n$ 重伯努利试验}。
    \item 若一个试验由可列个相互独立且同分布的子试验组成，则称其为\textbf{独立重复试验}。
\end{enumerate}
\end{mydefinition}
\origin{原笔记第1.4节，定义“独立试验、伯努利试验与独立重复试验”；源码第1853行。}

\subsection{例题：伯努利试验中的常用计算}\label{note:c1:e27}
\paragraph{题目}
进行独立重复伯努利试验，每次成功的概率为 $p$，失败的概率为 $q=1-p$。
其中 $0<p<1$。求：
\begin{enumerate}
\item 前 $n$ 次中恰好成功 $k$ 次的概率（$0\leq k\leq n$）。
\item 第一次成功发生在第 $k$ 次的概率（$k\geq1$）。
\item 第 $r$ 次成功发生在第 $k$ 次的概率（$1\leq r\leq k$）。
\item 永远不成功及最终至少成功一次的概率。
\item 第 $n$ 次成功出现在第 $m$ 次失败之前的概率（$m,n\geq1$），分别按第 $n$ 次成功前的失败次数、第 $m$ 次失败前的成功次数、前 $m+n-1$ 次成功总数分类，写出三种等价形式。
\end{enumerate}

\basis{原笔记第1.4节，例“伯努利试验中的常用计算”；源码\texttt{chapters/chapter1.tex}第1863行。}
\begin{proof}[解答]
\begin{enumerate}
\item 指定某 $k$ 次成功、其余 $n-k$ 次失败，概率为 $p^kq^{n-k}$。成功位置有 $\binom nk$ 种，各种情况互斥，故概率为
\[\binom nk p^kq^{n-k}.\]
\item 前 $k-1$ 次必须失败，第 $k$ 次成功，由独立性得
\[q^{k-1}p.\]
\item 前 $k-1$ 次中须恰好有 $r-1$ 次成功，第 $k$ 次还须成功，故概率为
\[\binom{k-1}{r-1}p^{r-1}q^{k-r}p=\binom{k-1}{r-1}p^rq^{k-r}.\]
\item 令 $F_n$ 表示前 $n$ 次全部失败，则
\[F_n\downarrow\bigcap_{n=1}^\infty F_n=\{\text{永远不成功}\}.\]
由概率的上连续性，
\[P(\text{永远不成功})=\lim_{n\to\infty}P(F_n)=\lim_{n\to\infty}q^n=0.\]
所以最终至少成功一次的概率为 $1$。
\item 记所求事件为 $E$。若第 $n$ 次成功前恰有 $k$ 次失败，则 $0\leq k\leq m-1$；前 $n+k-1$ 次有 $n-1$ 次成功，下一次成功。因此
\[P(E)=\sum_{k=0}^{m-1}\binom{k+n-1}{n-1}p^nq^k.\]
若按第 $m$ 次失败前的成功次数 $k$ 分类，则 $E$ 等价于 $k\geq n$。第 $m$ 次失败发生在第 $k+m$ 次，故
\[P(E)=\sum_{k=n}^{\infty}\binom{k+m-1}{k}p^kq^m.\]
将成功与失败互换，上一项说明以后至少还会失败一次的概率为 $1$。每次失败后的剩余试验仍为同样的独立重复试验，重复此论证并作有限次归纳，可知第 $m$ 次失败以概率 $1$ 最终发生。因此上述分类遗漏的事件只有零概率。

也可考察前 $m+n-1$ 次试验：$E$ 发生当且仅当其中至少有 $n$ 次成功。若至少有 $n$ 次成功，则失败至多 $m-1$ 次，第 $n$ 次成功已发生而第 $m$ 次失败尚未发生；反之，若成功至多 $n-1$ 次，则已有至少 $m$ 次失败，第 $m$ 次失败先发生。因此
\[P(E)=\sum_{k=n}^{m+n-1}\binom{m+n-1}{k}p^kq^{m+n-1-k}.\]
\end{enumerate}
\end{proof}

\subsection{例题：分赌注问题}\label{note:c1:e28}
\paragraph{题目}
甲、乙约定先胜 $t$ 局者获得全部赌注。比赛在甲已胜 $r$ 局、乙已胜 $s$ 局时中止，其中 $r<t$、$s<t$。以后每局甲获胜的概率为 $p$，且各局相互独立。求甲、乙最终获胜的概率，并说明中止时应如何分配全部赌注。

\basis{原笔记第1.4节，例“分赌注问题”；源码\texttt{chapters/chapter1.tex}第1908行。}
\begin{proof}[解答]
有
\[
\begin{aligned}
P(\text{甲最终先胜 }t\text{ 局})
&=P(t-r\text{ 次成功发生在 }t-s\text{ 次失败之前})\\
&=\sum_{k=0}^{t-s-1}
\binom{k+t-r-1}{t-r-1}p^{t-r}(1-p)^k.
\end{aligned}
\]
合理的分配方式是按甲、乙最终先胜 $t$ 局的概率分配赌注。记上式结果为 $u$，则乙最终先胜 $t$ 局的概率为 $1-u$，甲、乙应分别取得赌注的 $u$ 与 $1-u$ 部分。
\end{proof}

\subsection{例题：\texorpdfstring{事件 $A$ 先于事件 $B$ 发生}{事件 A 先于事件 B 发生}}\label{note:c1:e29}
\paragraph{题目}
进行独立重复试验。每次试验中，互斥事件 $A,B$ 发生的概率分别为 $p_A,p_B$，并设 $p_A+p_B>0$。求 $A$ 先于 $B$ 发生的概率，分别用按首次出现时刻求和和按第一次试验分类两种方法计算。

\basis{原笔记第1.4节，例“事件 $A$ 先于事件 $B$ 发生”；源码\texttt{chapters/chapter1.tex}第1921行。}
\begin{proof}[解答]
$A$ 首次出现且早于 $B$ 发生在第 $n$ 次，当且仅当前 $n-1$ 次二者均未发生，第 $n$ 次发生 $A$。这些事件两两互斥，故

\[
\begin{aligned}
P(A\text{ 先于 }B\text{ 发生})
&=\sum_{n=1}^{\infty}(1-p_A-p_B)^{n-1}p_A\\
&=\frac{p_A}{p_A+p_B}.
\end{aligned}
\]
也可以令所求概率为 $x$。按第一次试验分类，
\[
x=p_A+(1-p_A-p_B)x,
\]
解得同一结果。
\end{proof}

\subsection{例题：由独立性得到欧拉乘积公式}\label{note:c1:e30}
\paragraph{题目}
设 $s>1$，在正整数集合 $\mathbb Z_{>0}$ 上定义
\[
P(A)=\frac{\displaystyle\sum_{k\in A}k^{-s}}
{\displaystyle\sum_{k=1}^{\infty}k^{-s}}.
\]
将所有素数依次记为 $p_1,p_2,\ldots$，令
\[
A_i=\{mp_i:m\in\mathbb Z_{>0}\}.
\]
\begin{enumerate}
\item 证明事件列 $A_1,A_2,\ldots$ 相互独立。
\item 利用概率的连续性证明欧拉乘积公式
\[\sum_{k=1}^{\infty}\frac1{k^s}=\prod_{i=1}^{\infty}\left(1-\frac1{p_i^s}\right)^{-1}.\]
\end{enumerate}

\basis{原笔记第1.4节，例“由独立性得到欧拉乘积公式”；源码\texttt{chapters/chapter1.tex}第1937行。}
\begin{proof}
\begin{enumerate}
\item 由定义，
\[P(A_i)=\frac{\sum_{m=1}^\infty(mp_i)^{-s}}{\sum_{k=1}^\infty k^{-s}}=p_i^{-s}.\]
任取不同的 $i_1,\ldots,i_r$，有
\[
A_{i_1}\cap\cdots\cap A_{i_r}
=\{mp_{i_1}\cdots p_{i_r}:m\in\mathbb Z_{>0}\},
\]
所以
\[
P(A_{i_1}\cap\cdots\cap A_{i_r})
=\frac1{p_{i_1}^s\cdots p_{i_r}^s}
=\prod_{j=1}^{r}P(A_{i_j}).
\]
因此 $A_1,A_2,\ldots$ 相互独立。
\item 记 $C_n=\bigcap_{i=1}^n A_i^{\mathrm c}$，则 $C_n$ 递减。对每个 $n$，由有限个独立事件的补事件仍相互独立，
\[P(C_n)=\prod_{i=1}^n(1-p_i^{-s}).\]
又因为每个大于 $1$ 的正整数至少有一个素因子，
\[
\{1\}=\bigcap_{i=1}^{\infty}A_i^{\mathrm c}.
\]
由概率的上连续性和独立性，
\[
P(\{1\})
=\prod_{i=1}^{\infty}\left(1-\frac1{p_i^s}\right).
\]
另一方面，
\[
P(\{1\})=\left(\sum_{k=1}^{\infty}\frac1{k^s}\right)^{-1}.
\]
故得到欧拉乘积公式
\[
\sum_{k=1}^{\infty}\frac1{k^s}
=\prod_{i=1}^{\infty}\left(1-\frac1{p_i^s}\right)^{-1}.
\]
\end{enumerate}
\end{proof}

\clearpage\section{关于试验独立性的补充}\label{k:product}

设子试验 $1$ 和子试验 $2$ 的概率空间分别为
\[
(\Omega_1,\mathscr F_1,P_1),
\qquad
(\Omega_2,\mathscr F_2,P_2).
\]
两个子试验组成的大试验以
\[
\Omega=\Omega_1\times\Omega_2
\]
为样本空间。子试验 $1$ 中的事件 $A\in\mathscr F_1$ 在大样本空间中表示为 $A\times\Omega_2$；子试验 $2$ 中的事件 $B\in\mathscr F_2$ 表示为 $\Omega_1\times B$。两者同时发生对应
\[
(A\times\Omega_2)\cap(\Omega_1\times B)=A\times B.
\]

\subsection{例题：掷骰子与独立取球}\label{note:c1:e31}
\paragraph{题目}
子试验 $1$ 为掷一枚均匀骰子，$\Omega_1=\{1,2,3,4,5,6\}$；子试验 $2$ 为从红、白、蓝三个球中等概率取一个，$\Omega_2=\{\text{红},\text{白},\text{蓝}\}$。两个子试验相互独立。
\begin{enumerate}
\item 求每个联合结果 $(i,c)\in\Omega_1\times\Omega_2$ 的概率。
\item 对任意 $A\subseteq\Omega_1$、$B\subseteq\Omega_2$，验证 $P(A\times B)=P_1(A)P_2(B)$。
\end{enumerate}

\basis{原笔记第1.5节，例“掷骰子与独立取球”；源码\texttt{chapters/chapter1.tex}第1999行。}
\begin{proof}[解答]
\begin{enumerate}
\item 每个联合结果 $(i,c)\in\Omega_1\times\Omega_2$ 的概率为
\[
P(\{(i,c)\})=\frac16\cdot\frac13=\frac1{18}.
\]
\item $A\times B$ 中共有 $n(A)n(B)$ 个联合结果，每个结果的概率为 $1/18$，故
\[P(A\times B)=\frac{n(A)n(B)}{18}=\frac{n(A)}6\frac{n(B)}3=P_1(A)P_2(B).\]
\end{enumerate}
\end{proof}

\begin{mydefinition}{两个独立子试验的乘积空间表述}{two-independent-experiments-product}
在 $\Omega_1\times\Omega_2$ 上取乘积 $\sigma$ 域
\[
\mathscr F_1\otimes\mathscr F_2
=\sigma\bigl(\{A\times B:A\in\mathscr F_1,\ B\in\mathscr F_2\}\bigr).
\]
若其上的概率 $P$ 对任意 $A\in\mathscr F_1$、$B\in\mathscr F_2$ 都满足
\[
P(A\times B)=P_1(A)P_2(B),
\]
则称这两个子试验相互独立，并将 $P$ 记为 $P_1\otimes P_2$。
\end{mydefinition}
\origin{原笔记第1.5节，定义“两个独立子试验的乘积空间表述”；源码第2010行。}

\begin{mydefinition}{有限个独立子试验}{finite-independent-experiments-product}
设第 $i$ 个子试验的概率空间为 $(\Omega_i,\mathscr F_i,P_i)$，$i=1,\ldots,n$。在
\[
\Omega=\prod_{i=1}^{n}\Omega_i,
\qquad
\mathscr F=\bigotimes_{i=1}^{n}\mathscr F_i
\]
上，若概率 $P$ 对任意 $A_i\in\mathscr F_i$ 都满足
\[
P(A_1\times\cdots\times A_n)
=\prod_{i=1}^{n}P_i(A_i),
\]
则称这 $n$ 个子试验相互独立，并记
\[
P=P_1\otimes\cdots\otimes P_n.
\]
若各个概率空间完全相同，则它们组成 $n$ 重独立试验。
\end{mydefinition}
\origin{原笔记第1.5节，定义“有限个独立子试验”；源码第2034行。}

\subsection{例题：掷骰子后按点数取球}\label{note:c1:e32}
\paragraph{题目}
先掷一枚均匀骰子。若点数为偶数，则从红、白两球中等概率取一个；若点数为奇数，则取蓝球。用
\[
p(\omega_1,\{\omega_2\})
=P(\text{第二步取出 }\omega_2\mid\text{第一步点数为 }\omega_1)
\]
表示转移概率。其中 $\Omega_1=\{1,2,3,4,5,6\}$，$\Omega_2=\{\text{红},\text{白},\text{蓝}\}$。
\begin{enumerate}
\item 列出各点数条件下取出各颜色球的转移概率。
\item 对任意 $A\subseteq\Omega_1$、$B\subseteq\Omega_2$，写出联合概率 $P(A\times B)$ 的表达式。
\item 求取出红球、骰子点数为 $2$、以及点数为 $2$ 且取出红球的概率，并判断两个子试验是否独立。
\end{enumerate}

\basis{原笔记第1.5节，例“掷骰子后按点数取球”；源码\texttt{chapters/chapter1.tex}第2055行。}
\begin{proof}[解答]
\begin{enumerate}
\item 若 $\omega_1=2,4,6$，则
\[
p(\omega_1,\{\text{红}\})=p(\omega_1,\{\text{白}\})=\frac12,
\qquad
p(\omega_1,\{\text{蓝}\})=0;
\]
若 $\omega_1=1,3,5$，则
\[
p(\omega_1,\{\text{红}\})=p(\omega_1,\{\text{白}\})=0,
\qquad
p(\omega_1,\{\text{蓝}\})=1.
\]
\item 由全概率公式，对事件 $A\subseteq\Omega_1$ 和颜色集合 $B\subseteq\Omega_2$，联合概率可写为
\[
P(A\times B)
=\sum_{\omega_1\in A}P_1(\{\omega_1\})
\sum_{\omega_2\in B}p(\omega_1,\{\omega_2\}).
\]
\item 有
\[
P(\text{取出红球})=\frac14,
\qquad
P(\text{骰子点数为 }2)=\frac16,
\]
而
\[
P(\text{点数为 }2\text{ 且取出红球})=\frac1{12}
\ne\frac16\cdot\frac14.
\]
所以两个子试验不独立。
\end{enumerate}
\end{proof}
